% Discussão e Ameaças à Validade (seções escritas pelos três).
% Números de resultados/, validacao/kappa.csv e dados/funil.csv.

\section{Discussão}

\textbf{As hipóteses sobre a taxa de falha erraram nas duas direções.} Esperávamos o CFR (a)
entre 10\% e 25\%, e ele ficou em 4,9\%. Os \textit{runs} contados são de \textit{push} no
\textit{default branch}, ou seja, de mudanças que já passaram por um \textit{pull request}
com CI; a métrica mede a falha do que já foi aprovado, que em projetos populares com revisão
obrigatória é baixa. Para o CFR (b) esperávamos valores iguais ou menores, e a mediana foi
quatro vezes maior. Parte disso vem do instrumento: a heurística adotada tem mais falsos
positivos (57) que falsos negativos (30) na amostra-ouro, então o valor de 0,200 tende a
superestimar o real. A única hipótese confirmada na RQ~03 foi a de que as variantes
concordariam pouco.

\textbf{As duas variantes do CFR medem coisas diferentes.} Com $\rho = 0{,}107$ e IC que
encosta no zero, um projeto bem avaliado por uma variante pode ser mal avaliado pela outra.
Recalculando a classificação geral com o CFR (b) no lugar do (a), 72 dos 350 repositórios
(20,6\%) mudam de categoria e a faixa \textit{Elite} cai de 25 para 8. No \textit{lead time} foi
diferente ($\rho = 0{,}744$ e medianas próximas), porque as duas variantes medem o mesmo evento
de formas diferentes, enquanto as do CFR medem eventos distintos: falha de \textit{pipeline}
antes da entrega e correção depois dela.

\textbf{A escala DORA e as releases open-source.} A frequência é a dimensão em que a amostra vai
pior, com 74,3\% dos repositórios abaixo de uma \textit{release} por semana. Isso diz mais sobre o
\textit{proxy} do que sobre a capacidade de entrega: uma \textit{release} é um marco de
versionamento, muitas vezes mensal ou trimestral, enquanto o \textit{deploy} do
DORA~\cite{dora2024} é a chegada de uma mudança à produção, e o patamar \textit{Elite} de sete por
semana pressupõe entrega contínua em ambiente corporativo~\cite{forsgren2018}. Como a categoria
geral é a mediana de quatro notas, ela acaba definida pelo par CFR (a) e recuperação, as
dimensões em que a amostra vai melhor.

\textbf{Velocidade e estabilidade (RQ 05).} As duas previsões se confirmaram na direção. Na
variante (a), a correlação é positiva mas fraca, e o IC exclui o zero por pouco: publicar com
mais frequência quase não vem acompanhado de mais falhas de CI, o que é compatível com a tese do
DORA de que velocidade e estabilidade não são um \textit{trade-off}. Na variante (b), a
correlação forte é em boa parte artefato da métrica, porque quem publica toda semana tem muito
mais chances de ter uma \textit{release} corretiva em até 7 dias do que quem publica a cada
trimestre. Com apenas a variante (b), a conclusão seria o oposto do que o DORA sustenta.

\textbf{O que explica as diferenças (RQ 06).} A idade é o fator mais associado ao desempenho:
projetos mais antigos acumulam mais tempo entre o \textit{commit} mais antigo e a publicação
($\varepsilon^2$ de 0,288 e 0,251 no \textit{lead time}) e publicam menos. A popularidade
quase não explica nada: o número de estrelas só aparece no \textit{lead time} (b), com o menor
efeito ($\varepsilon^2 = 0{,}049$).

\textbf{As conclusões são robustas? (RQ 07 e RQ 08).} Depende da RQ. A do \textit{lead time}
(RQ~02) se mantém nas duas variantes. As da taxa de falha (RQ~03 e 05) e a classificação geral
não se mantêm: trocar a variante do CFR muda a categoria de 20,6\% dos repositórios, e trocar a
unidade de \textit{deploy} muda entre 25\% e 33\%, e nenhum desses ICs desce abaixo de 21\%. O
\textit{rework rate} reforça a leitura do CFR (b): metade das \textit{releases} corrige algo, mas
só uma em cinco corrige algo publicado na semana anterior, então a maior parte da correção é
manutenção de rotina.

\section{Ameaças à Validade}

Seguimos as quatro categorias de Wohlin et al.~\cite{wohlin2012}.

\textbf{Construto.} \textit{Release} não é \textit{deploy}, e falha de CI não é falha em
produção: um \textit{run} que falha costuma ser um erro barrado antes de chegar ao usuário. A
validação manual sustentou o primeiro \textit{proxy} (as \textit{releases} representam entregas
reais em 59 dos 60 repositórios da amostra-ouro), mas a heurística de \textit{release} corretiva
não atingiu o F1 de 0,70 (0,65 na v4), e esse erro passa para o CFR (b), a RQ~08 e a C2 da
RQ~07. Os próprios avaliadores divergiram em 26 de 299 \textit{releases}. Rebase e
\textit{squash merge} reescrevem o \texttt{commit.author.date} e afetam o \textit{lead time}
em direção desconhecida.

\textbf{Interna.} As análises são correlacionais. A idade se correlaciona com a frequência
($\rho = -0{,}306$), assim como as estrelas ($\rho = 0{,}162$) e os contribuidores
($\rho = 0{,}119$), e esses fatores não são controlados nas RQ~01 a 04. O filtro de inclusão
(5 \textit{releases} e 50 \textit{runs}) deixa só projetos que já versionam e automatizam com
regularidade, e o descarte dos seis repositórios com mais de 1.000 \textit{runs} por dia tira
da amostra os de maior atividade de CI.

\textbf{Externa.} A amostra não é aleatória: entraram os 1.000 mais estrelados de cada faixa, e
os menos estrelados de cada faixa nunca foram observados. Os resultados valem para repositórios
populares com GitHub Actions, e não se transferem para equipes corporativas, que são o contexto
das faixas DORA. A janela é única, de 12 meses, e não separa característica estável de efeito
do período.

\textbf{Conclusão.} Usamos mediana, IQR e testes não paramétricos por causa da assimetria, e a
correção de Holm nos 30 testes da RQ~06. O tipo do projeto só existe para 60 repositórios, e os
testes sobre ele têm pouco poder. Os ICs por \textit{bootstrap} mostram que as medianas e as
proporções de mudança de categoria são estáveis, mas que a correlação entre as variantes do CFR
não se distingue com segurança de zero.
